\begin{algorithm}[H]
\caption{Evolved selection: suppress clones and explore quality-gated niches}
\label{alg:evolved-selection}
\io{Input}{Population $P$; tournament size $k$; geometric parameter $p$; maximum
complexity $L$; global normalized complexity frequencies $f_c$; parsimony scale $\lambda$.}
\begin{algorithmic}[1]
\State $S\gets$ uniformly sample $\min(k,|P|)$ members without replacement
\For{$i\in S$}
  \State $c_i\gets\fn{Complexity}(i)$; $\ell_i\gets\fn{Loss}(i)$
  \State $s_i\gets\fn{Cost}(i)\exp(\lambda f_{c_i})$ if frequency weighting is enabled,
  \Statex \hspace{\algorithmicindent}otherwise $s_i\gets\fn{Cost}(i)$; use $f_c=0$ outside $1\le c\le L$
\EndFor
\State $\pi\gets\fn{Argsort}(s)$; $\widetilde s\gets s$
\For{$a=2,\ldots,|S|$}
  \State $i\gets\pi_a$
  \If{some $j\in\{\pi_1,\ldots,\pi_{a-1}\}$ satisfies
    $c_i=c_j$ and $|\ell_i-\ell_j|\le10^{-6}\max(|\ell_i|,|\ell_j|,\epsilon)$}
    \State $\widetilde s_i\gets\begin{cases}
      1.35s_i,&s_i>0,\\s_i/1.35,&s_i<0,\\0,&s_i=0
    \end{cases}$ \quad\Comment{One sign-safe penalty per clone}
  \EndIf
\EndFor
\State $\rho\gets\fn{Argsort}(\widetilde s)$
\If{$\Unif(0,1)<0.10$}
  \State $H\gets$ first member of each valid complexity in $\rho$ order
  \If{$H\ne\varnothing$}
    \State $H_q\gets$ first $\lceil |H|/2\rceil$ members of $H$ \quad\Comment{Keep the better half}
    \State \Return $\argmin_{i\in H_q}(f_{c_i},\widetilde s_i)$ \quad\Comment{Rarest niche; quality breaks ties}
  \EndIf
\EndIf
\State Draw rank $r\in\{1,\ldots,|S|\}$ with $\Pr(r)\propto p(1-p)^{r-1}$
\State \Return $S_{\rho_r}$ \quad\Comment{For $p=1$, choose the best adjusted score}
\end{algorithmic}
\algnotes{Clone detection uses complexity and nearly equal loss as a phenotype proxy;
it does not compare expression strings. All earlier members in the \emph{original} ranking
can establish a clone match, including earlier clones. Niche rarity uses global search
frequencies, even when frequency weighting of tournament costs is disabled.
If no valid niche exists, selection falls back to the geometric tournament.}
\end{algorithm}
